\documentclass[11pt]{article}

\usepackage{amsmath,amssymb,amsthm}
\usepackage[margin=1in]{geometry}
\usepackage{hyperref}

\newtheorem{theorem}{Theorem}
\newtheorem{lemma}{Lemma}
\newtheorem{remark}{Remark}

\DeclareMathOperator{\SubSum}{SubSum}

\title{A Fixed Covering Certificate for Subset Sums of Powers of \(3,4,7\)}
\author{[Author Name]}
\date{}

\begin{document}

\maketitle

\begin{abstract}
Let
\[
B=\{3^a:a\ge1\}\cup\{4^b:b\ge1\}\cup\{7^c:c\ge1\}.
\]
We prove that every integer \(n\ge582\) is representable as a subset sum of distinct
elements of \(B\). The proof combines a finite subset-sum interval certificate with
an interval-overlap propagation lemma. The principal finite check is a weighted
gap lemma for later powers of \(3,4,7\), verified by a replayable modular certificate.
The result is a fixed-case covering theorem only; no optimality of the endpoint is
claimed.
\end{abstract}

\section{Introduction}

Questions about which integers are representable as subset sums of a prescribed
sequence are a classical elementary theme.  For nondecreasing sequences beginning
with \(1\), Brown's criterion gives a well-known interval-overlap condition for
completeness; see \cite{brown1961}.  The present note treats a narrower eventual
covering problem for a sparse sequence that does not begin with \(1\).

Let
\[
B=\{3^a:a\ge1\}\cup\{4^b:b\ge1\}\cup\{7^c:c\ge1\}.
\]
We prove that every integer at least \(582\) is a subset sum of distinct elements of
\(B\).  The point of recording the result is not the size of the endpoint, and no
minimality assertion is made.  Rather, the proof gives a compact fixed-case
certificate: a finite interval of coverage is verified directly, and a finite
modular certificate verifies that the overlap condition continues for every later
power in the sequence.

The proof begins with a finite interval certificate
\[
[582,4460]\subset \SubSum(\{3,4,7,9,16,27,49,64,81,243,256,343,729,1024,2187\}),
\]
then uses an interval-overlap lemma to propagate coverage indefinitely.  The
propagation step is reduced to a weighted gap lemma for the next powers of
\(3,4,7\).  That lemma is verified by a finite modular certificate supplied with the
paper.

\section{Definitions}

Let \(B\) be ordered increasingly as
\[
b_1<b_2<b_3<\cdots .
\]
For a finite set \(S\) of integers, define
\[
\SubSum(S)=\left\{\sum_{x\in T}x:T\subseteq S\right\}.
\]
We write
\[
A_1(\{3,4,7\})=\SubSum(B).
\]
For \(t\in B\), define
\[
S(t)=\sum_{\substack{p\in B\\ p<t}}p.
\]
For \(t>2187\), define the next-power gaps
\[
\Delta_3(t)=\min\{3^a:3^a>t\}-t,
\]
\[
\Delta_4(t)=\min\{4^b:4^b>t\}-t,
\]
and
\[
\Delta_7(t)=\min\{7^c:7^c>t\}-t.
\]

\section{The fixed covering theorem}

\begin{theorem}
Let
\[
B=\{3^a:a\ge1\}\cup\{4^b:b\ge1\}\cup\{7^c:c\ge1\}.
\]
Then
\[
[582,\infty)\subset A_1(\{3,4,7\}).
\]
Equivalently, every integer \(n\ge582\) is representable as a subset sum of
distinct elements of \(B\).
\end{theorem}

\section{The initial interval certificate}

The first fifteen elements of \(B\) are
\[
3,4,7,9,16,27,49,64,81,243,256,343,729,1024,2187.
\]
A finite subset-sum enumeration verifies
\[
[582,4460]\subset
\SubSum(\{3,4,7,9,16,27,49,64,81,243,256,343,729,1024,2187\}).
\]
Set
\[
M=582,\qquad N=4460.
\]
The sum of the first fifteen listed terms is \(5042\), so \(N=5042-M\).

\section{Interval-overlap propagation}

\begin{lemma}[Interval-overlap propagation]
Let
\[
[M,N]\subset \SubSum(S)
\]
for a finite set \(S\) of distinct allowed terms. Let \(p\notin S\) be a new allowed
term. If
\[
p\le N-M+1,
\]
then
\[
[M,N+p]\subset \SubSum(S\cup\{p\}).
\]
\end{lemma}

\begin{proof}
The interval \([M,N]\) is already covered by subset sums from \(S\). Adding \(p\)
to each representation covers the interval
\[
[M+p,N+p].
\]
The condition \(p\le N-M+1\) is equivalent to
\[
M+p\le N+1.
\]
Hence \([M,N]\) and \([M+p,N+p]\) overlap or touch. Their union therefore covers
\([M,N+p]\), proving the claim.
\end{proof}

Suppose the current covered interval has the form \([M,N]\), where \(M=582\), and
the next unused power is \(t\in B\). In the propagation below, \(N=S(t)-M\), so the
condition \(t\le N-M+1\) is equivalent to
\[
S(t)-t\ge 2M-1=1163.
\]
It remains to prove this inequality for every \(t>2187\) in \(B\).

\section{Weighted next-gap identities}

\begin{lemma}[Weighted next-gap identities]
For every future power \(t>2187\) generated by \(3,4,7\), the following identities
hold.

If \(t=3^r\), \(r\ge8\), then
\[
S(t)-t-1163=\frac{2\Delta_4(t)+\Delta_7(t)-7002}{6}.
\]

If \(t=4^s\), \(s\ge6\), then
\[
S(t)-t-1163=\frac{3\Delta_3(t)+\Delta_7(t)-7002}{6}.
\]

If \(t=7^q\), \(q\ge4\), then
\[
S(t)-t-1163=\frac{3\Delta_3(t)+2\Delta_4(t)-7002}{6}.
\]
\end{lemma}

\begin{proof}
We prove the first identity; the other two follow by the same calculation. Let
\(t=3^r\), and let
\[
4^a=3^r+\Delta_4(t),\qquad 7^b=3^r+\Delta_7(t)
\]
be the least larger powers of \(4\) and \(7\), respectively. Since the powers below
\(3^r\) are finite geometric sums on the three bases,
\[
S(3^r)=\frac{3^r-3}{2}+\frac{4^a-4}{3}+\frac{7^b-7}{6}.
\]
Substituting \(4^a=3^r+\Delta_4(t)\) and
\(7^b=3^r+\Delta_7(t)\) gives
\[
S(3^r)-3^r-1163
=
\frac{2\Delta_4(t)+\Delta_7(t)-7002}{6}.
\]
The identities for \(t=4^s\) and \(t=7^q\) are obtained similarly, using the least
larger powers on the two competing bases.
\end{proof}

\section{The weighted gap lemma}

\begin{lemma}[Weighted gap lemma]
For every future power \(t>2187\) generated by \(3,4,7\), the corresponding
weighted gap inequality holds:
\[
t=3^r,\ r\ge8
\quad\Longrightarrow\quad
2\Delta_4(t)+\Delta_7(t)\ge7002,
\]
\[
t=4^s,\ s\ge6
\quad\Longrightarrow\quad
3\Delta_3(t)+\Delta_7(t)\ge7002,
\]
and
\[
t=7^q,\ q\ge4
\quad\Longrightarrow\quad
3\Delta_3(t)+2\Delta_4(t)\ge7002.
\]
\end{lemma}

\begin{proof}
The proof is by a finite modular certificate. A failure of the first inequality
would give a bounded full paired system
\[
4^a-3^r=u,\qquad 7^b-3^r=v,\qquad 2u+v\le7001,\qquad r\ge8,
\]
with \(1\le u\le3500\) and \(1\le v\le7000\).

A failure of the second inequality would give
\[
3^a-4^s=u,\qquad 7^b-4^s=v,\qquad 3u+v\le7001,\qquad s\ge6,
\]
with \(1\le u\le2333\) and \(1\le v\le7000\).

A failure of the third inequality would give
\[
3^a-7^q=u,\qquad 4^b-7^q=v,\qquad 3u+2v\le7001,\qquad q\ge4,
\]
with \(1\le u\le2333\) and \(1\le v\le3500\).

The residue filters are finite. In the first case, \(r\ge8\) implies
\[
3^r\equiv0\pmod{6561},
\]
so every failure must satisfy
\[
4^a\equiv u\pmod{6561},\qquad 7^b\equiv v\pmod{6561}.
\]
In the second case, \(s\ge6\) implies
\[
4^s\equiv0\pmod{4096},
\]
so every failure must satisfy
\[
3^a\equiv u\pmod{4096},\qquad 7^b\equiv v\pmod{4096}.
\]
In the third case, \(q\ge4\) implies
\[
7^q\equiv0\pmod{2401},
\]
so every failure must satisfy
\[
3^a\equiv u\pmod{2401},\qquad 4^b\equiv v\pmod{2401}.
\]

The accompanying replayable certificate constructs all residue-admissible systems
and eliminates them by full paired modular tests. The initial residue-admissible
counts are
\[
1,361,889,\qquad 255,938,\qquad 1,501,666
\]
in the three cases, respectively. In each case, the final number of surviving
systems is zero. The certificate checks the full paired systems displayed above,
not merely the resulting difference equations. Therefore every possible failure
system is excluded, and the weighted gap inequalities hold.
\end{proof}

\section{Proof of the theorem}

The finite interval certificate gives
\[
[582,4460]\subset \SubSum(b_1,\ldots,b_{15}),
\]
where
\[
b_1,\ldots,b_{15}
=
3,4,7,9,16,27,49,64,81,243,256,343,729,1024,2187.
\]
Let \(t=b_j>2187\) be any later power in the ordered sequence \(B\). By the
weighted gap lemma and the weighted next-gap identities,
\[
S(t)-t\ge1163.
\]
As observed above, this is exactly the interval-overlap condition for adjoining
the next power \(t\). Hence the interval-overlap lemma applies successively to
every later power. Starting from \([582,4460]\), the covered interval propagates
without bound. Therefore
\[
[582,\infty)\subset A_1(\{3,4,7\}).
\]
This proves the theorem.

\section{Computational certificate}

The supplementary files contain a replayable computational certificate.  The
certificate checks two finite claims.

First, it checks the initial subset-sum interval certificate
\[
[582,4460]\subset
\SubSum(\{3,4,7,9,16,27,49,64,81,243,256,343,729,1024,2187\}).
\]
This is finite because only \(2^{15}\) subset sums are involved.

Second, it checks the weighted gap lemma.  Each possible failure of the lemma gives
one of the bounded full paired systems described in the proof.  For example, in
Case A a failure is represented by integers \(u,v\) satisfying
\(1\le u\le3500\), \(1\le v\le7000\), and \(2u+v\le7001\), together with the full
paired system
\[
4^a-3^r=u,\qquad 7^b-3^r=v,\qquad r\ge8.
\]
The other two cases are analogous.  Thus each case begins with finitely many
bounded pairs \((u,v)\).

For each bounded pair, the verifier first applies the stabilized residue filter
modulo \(6561\), \(4096\), or \(2401\), depending on the case.  This produces a
finite list of residue-admissible systems, with forced exponent classes modulo the
corresponding multiplicative orders.  The verifier then applies the moduli listed
in the certificate tables.  A row is eliminated only when the relevant congruences
for the full paired system have no simultaneous solution modulo the selected
modulus.

This last point is important: the verifier does not merely test the difference
equation obtained by subtracting the two displayed equations.  It checks the
paired congruence system involving the common baseline power.  The certificate
therefore excludes the full possible failure systems required by the weighted gap
lemma.

The certificate status is:

\begin{quote}
locally replayed; replayable computational certificate included.
\end{quote}

No independent third-party acceptance of the certificate is asserted here.

To replay the certificates, run:
\begin{verbatim}
python3 verify_initial_interval.py
python3 verify_weighted_gap_certificates.py
\end{verbatim}
The expected weighted-gap terminal summary is:
\begin{verbatim}
WEIGHTED GAP CERTIFICATE REPLAY — PASS
Case A: initial = 1,361,889; final survivors = 0
Case B: initial = 255,938; final survivors = 0
Case C: initial = 1,501,666; final survivors = 0
Total residue-admissible systems checked = 3,119,493
Weighted Gap Lemma: PASS
Global overlap inequality: PASS
Fixed theorem chain: PASS
FINAL STATUS: PASS — WEIGHTED_GAP_PROVED
\end{verbatim}

\section{Limitations}

This paper proves only the fixed case \(D=\{3,4,7\}\) with \(k=1\). It does not
claim a result for any other generating set. It does not claim a result for
\(k\ne1\). It does not claim that \(582\) is the least possible endpoint. It does
not claim a density theorem, an asymptotic theorem, or a classification theorem.


\begin{thebibliography}{9}

\bibitem{brown1961}
J. L. Brown, Jr.,
\emph{Notes on complete sequences of integers},
American Mathematical Monthly \textbf{68} (1961), 557--560.

\end{thebibliography}

\end{document}
